% Part: incompleteness
% Chapter: arithmetization-syntax
% Section: coding-terms

\documentclass[../../../include/open-logic-section]{subfiles}

\begin{document}

\olfileid{inc}{art}{trm}
\olsection{رمزگذاری ترم‌ها}

\begin{explain}
ترم صرفاً نوع معینی از دنبالهٔ نمادهاست: طبق قواعد تشکیل ترم‌ها به‌طور
استقرایی از ثابت‌ها و متغیرها ساخته می‌شود. چون دنباله‌های نمادها را
می‌توان با اعداد رمزگذاری کرد، با استفاده از طرحی برای رمزگذاری
نمادها همراه با روشی برای رمزگذاری دنباله‌های اعداد، نسبت‌دادن اعداد
گودل به ترم‌ها دشوار نیست. چالش اصلی نشان‌دادن این است که ویژگیِ عدد
گودلِ یک ترم خوش‌ساخت بودن محاسبه‌پذیر، و در واقع بازگشتی اولیه، است.
\end{explain}

!!^{variable}s و !!{constant}s ساده‌ترین ترم‌ها هستند و آزمودن اینکه آیا
$x$ عدد گودل چنین ترمی است آسان است: $\fn{Var}(x)$ برقرار است اگر
$x$ برای بعضی~$i$ برابر~$\Gn{\Obj v_i}$ باشد. به بیان دیگر، $x$
دنباله‌ای به طول~$1$ است و تنها عضو آن~$(x)_0$ رمزِ !!{variable}ی چون
$\Obj v_i$ است؛ یعنی $x$ برای بعضی~$i$ برابر
$\tuple{\tuple{1,i}}$ است. به همین ترتیب، $\fn{Const}(x)$ برقرار است
اگر $x$ برای بعضی~$i$ برابر~$\Gn{\Obj c_i}$ باشد. هر دو رابطه
بازگشتی اولیه‌اند، زیرا اگر چنین~$i$ای وجود داشته باشد باید $<x$
باشد:
\begin{align*}
  \fn{Var}(x) & \defiff \bexists{i<x}{x = \tuple{\tuple{1, i}}}\\
  \fn{Const}(x) & \defiff \bexists{i<x}{x = \tuple{\tuple{2, i}}}
\end{align*}

\begin{prop}
\ollabel{prop:term-primrec}
روابط~$\fn{Term}(x)$ و~$\fn{ClTerm}(x)$ که به‌ترتیب اگر و تنها اگر
$x$ عدد گودل یک ترم یا یک ترم بسته باشد برقرارند، بازگشتی اولیه‌اند.
\end{prop}

\begin{proof}
دنبالهٔ نمادهای~$s$ ترم است اگر و تنها اگر دنباله‌ای از ترم‌ها چون
$s_0$،~\dots، $s_{k-1}=s$ وجود داشته باشد که ثبت کند ترم~$s$ چگونه
طبق قواعد تشکیل ترم‌ها از !!{constant}s و !!{variable}s ساخته شده است. یک
دنبالهٔ تشکیلِ فرضی برای آنکه از قواعد پیروی کند باید به ازای هر
$i<k$ یکی از شرایط زیر را داشته باشد:
\begin{enumerate}
\item $s_i$ !!a{variable}ی چون~$\Obj v_j$ باشد؛ یا
\item $s_i$ !!a{constant}ی چون~$\Obj c_j$ باشد؛ یا
\item $s_i$ با استفاده از !!{function} $n$موضعیِ~$\Obj f_j^n$ از~$n$
  ترم~$t_1$،~\dots، $t_n$ ساخته شده باشد که پیش از جایگاه~$i$ آمده‌اند.
\end{enumerate}
برای نشان‌دادن بازگشتی اولیه‌بودن رابطهٔ متناظر روی اعداد گودل، باید
این شرط را به‌طور بازگشتی اولیه، یعنی با توابع و روابط بازگشتی اولیه
و سورگذاری کراندار، بیان کنیم.

فرض کنید $y$ عددی باشد که دنبالهٔ~$s_0$،~\dots، $s_{k-1}$ را رمزگذاری
می‌کند؛ یعنی
$y=\tuple{\Gn{s_0},\dots,\Gn{s_{k-1}}}$. این عدد دنبالهٔ تشکیل ترمی
با عدد گودل~$x$ را رمزگذاری می‌کند اگر و تنها اگر برای هر~$i<k$:
\begin{enumerate}
\item $\fn{Var}((y)_i)$؛ یا
\item $\fn{Const}((y)_i)$؛ یا
\item $n$ای و عددی چون~$z=\tuple{z_1,\dots,z_n}$ وجود داشته باشد
  به‌گونه‌ای که هر~$z_l$ با بعضی~$(y)_{i'}$ برای~$i'<i$ برابر باشد و
\[
(y)_i = \Gn{\Obj f^n_j(} \concat \fn{flatten}(z) \concat \Gn{)},
\]
\end{enumerate}
و افزون بر این، $(y)_{k-1}=x$. (تابع~$\fn{flatten}(z)$ دنبالهٔ
$\tuple{\Gn{t_1},\dots,\Gn{t_n}}$ را به~$\Gn{t_1,\dots,t_n}$ تبدیل
می‌کند و بازگشتی اولیه است.)

نمایه‌های~$j$ و~$n$، اعداد گودل~$z_l$ برای ترم‌های~$t_l$، و رمز~$z$
برای دنبالهٔ~$\tuple{z_1,\dots,z_n}$ در بند (3) همگی از~$y$
کوچک‌ترند. در بالا می‌توانیم $k$ را با~$\len{y}$ جایگزین کنیم. پس
می‌توانیم «$y$ رمز دنبالهٔ تشکیل ترم دارای عدد گودل~$x$ است» را به
شیوه‌ای بیان کنیم که بازگشتی اولیه‌بودن این رابطه را نشان دهد.

اکنون فقط باید خود را قانع کنیم که کرانی بازگشتی اولیه برای~$y$ وجود
دارد. اما اگر $x$ عدد گودل ترمی باشد، باید دنبالهٔ تشکیلی با حداکثر
$\len{x}$ ترم داشته باشد (زیرا هر ترم در دنبالهٔ تشکیل~$s$ باید از
جایگاهی در~$s$ آغاز شود و هیچ دو زیرترمی نمی‌توانند از یک جایگاه آغاز
شوند). عدد گودل هر زیرترم~$s$ البته $\le x$ است. پس همیشه دنبالهٔ
تشکیلی با رمز~$\le p_{k-1}^{k(x+1)}$ وجود دارد که در آن
$k=\len{x}$.

برای~$\fn{ClTerm}$ کافی است بند مربوط به !!{variable}s را کنار بگذاریم.
\end{proof}

\begin{prob}
نشان دهید تابع~$\fn{flatten}(z)$ که دنبالهٔ
$\tuple{\Gn{t_1},\dots,\Gn{t_n}}$ را به~$\Gn{t_1,\dots,t_n}$ تبدیل
می‌کند، بازگشتی اولیه است.
\end{prob}

\begin{prop}\ollabel{prop:num-primrec}
  تابع~$\fn{num}(n)=\Gn{\num{n}}$ بازگشتی اولیه است.
\end{prop}

\begin{proof}
  $\fn{num}(n)$ را با بازگشت اولیه تعریف می‌کنیم:
  \begin{align*}
    \fn{num}(0) & = \Gn{\Obj 0}\\
    \fn{num}(n+1) & = \Gn{\prime(} \concat \fn{num}(n) \concat \Gn{)}.
  \end{align*}
\end{proof}

\end{document}
